% Analysis of the results. The numbers are macros from numbers.tex (make_numbers.py).
\newcommand{\rqOneText}{dsp-props detects \killFuncDsp{} of the \nFunctional{}
functional mutants and the ported suites \killFuncSop{}; neither raises an
alarm on an equivalent mutant, and of the ablations only generic PBT does
(Section~\ref{sec:rq2}). The two sets of detected faults are not nested:
\killSopOnly{} mutants are detected only by the suites, \killDspOnly{} only by
dsp-props, and together they detect \killUnion{} of them.

Every miss of dsp-props has the same cause: the fault shows only for operator
parameters that no declaration uses. A declaration fixes the coefficients of a
filter or the rates of a resampler, and dsp-props then varies the input signal.
Seven \code{lfilter} mutants break filters of order 0 or 1, or filters whose
numerator is shorter than the denominator, and the declared filters have orders
4 and 2 with numerators at least as long. The \code{upfirdn} survivors change the
length of the full output, which the declarations never see since they reach
\code{upfirdn} through \code{resample\_poly} and through a FIR filter truncated to
the input length, or ignore the last input sample, which every relation sees
identically on both of its sides. The suites miss the opposite kind of fault.
Their only comparisons with expected values on the FFT use length 30, which the
subject computes with Bluestein's algorithm; on powers of two they only check
that \code{ifft} inverts \code{fft}, which two mutants of the radix-2 twiddle
factors still satisfy. They also always pass the rates of \code{upfirdn}
explicitly, so four mutants of its default arguments survive.

dsp-props therefore does not detect more of these faults than suites refined
over many years; it detects other ones, without expected values, and its misses
are a matter of declarations rather than of relations: declaring a first-order
filter would have caught the \code{lfilter} survivors.}

\newcommand{\rqTwoText}{On the generated mutants the four randomised variants
detect almost the same faults (Table~\ref{tab:comparison}). The exceptions are
three FFT mutants that the variants with uniform inputs miss in some seeds,
while dsp-props detects them in every seed; two of them are the twiddle-factor
mutants that the ported suites miss. Generic PBT also raised a false alarm in
one of its \nSeeds{} seeds, on the correct FFT and on every FFT mutant
equivalent to it: Hypothesis drew a subnormal input of norm $10^{-323}$, for
which a purely relative tolerance is zero and a rounding error of one subnormal
fails Parseval's identity. The error model of dsp-props includes an absolute
floor, and neither dsp-props nor its other ablations raised an alarm. The
other mutants change the output by orders of magnitude, and almost any input
reveals them:
the median number of inputs before the first failure is \firstMedDsp{} for
dsp-props and \firstMedUni{} for uniform inputs, and the median counterexample
has \cexMedDsp{} and \cexMedUni{} samples. The generators and tolerances do
not change what is detected here, and the generators cost time, with a median of
\secMedDsp{}~s per mutant against \secMedUni{}~s for uniform inputs and
\secMedSop{}~s for a ported suite.

The hand-written faults, which include the accuracy faults, separate the
variants (Table~\ref{tab:hand}). Every variant detects every fault in every seed,
except the fixed tolerance, which misses one resampler fault in one seed, but
not equally fast. With the domain generators, the median failure appears at
input \handDspFirstMed{} (90th percentile \handDspFirstPnn{}); with uniform inputs
at input \handUniFirstMed{} (\handUniFirstPnn{}), and with generic PBT at input
\handGenFirstMed{} (\handGenFirstPnn{}). Counterexamples are smaller too: the 90th percentile of their length is
\handDspCexPnn{} samples for dsp-props, \handUniCexPnn{} with uniform inputs, and
\handFixCexPnn{} with a fixed tolerance. No variant raised a false alarm on
the correct references of this table. With 200 inputs per relation, these differences do not change the detection
rate. We expect them to matter with a smaller budget or on faults that need a
specific input, but we did not measure this. One hand-written fault, a filter
whose gain is wrong by 1\,\%, is detected only because its transfer function is
declared: it is still linear, time-invariant and causal.}

\newcommand{\conclusionResult}{On \nFunctional{} generated faults it detects
\killFuncDsp{}, against \killFuncSop{} for the ported test suites of NumPy and
SciPy, and the two are complementary: together they detect \killUnion. Its
misses come from operator parameters that no declaration exercised; the misses
of the suites, from code paths that their few expected values do not reach.}

\newcommand{\reportedUpstream}{We reported the empty-block faults to
SciPy\footnote{\url{https://github.com/scipy/scipy/issues/26406}} and
contributed a fix\footnote{\url{https://github.com/scipy/scipy/pull/26415}},
the dropped delay of \code{tf2sos} to
SciPy\footnote{\url{https://github.com/scipy/scipy/issues/26417}}, and the
clamping to torchaudio\footnote{\url{https://github.com/pytorch/audio/issues/4238}}.}
